% Minimal preamble for Open Logic Project sources (pinned commit 1e960beff9ed7835bf3e3f1335e21af3439cd107).
%
% OLP's own master preamble cannot run on the project host: TeX Live 2025 there lacks
% prooftree.sty, ntheorem.sty, extarrows.sty, mathdots.sty, isomath.sty and subfiles.sty.
% Following the project's own rule ("if the author preamble needs packages the toolchain lacks,
% do not invent macros -- use a minimal working preamble"), this file declares, verbatim in
% spirit and numbering behaviour, exactly the theorem-like environments OLP defines in
% sty/open-logic-book-envs.sty (thm/lem/prop/cor/defn/ex/prob/rem/note/case/axiom/conv),
% plus the shorthand accents OLP uses (`!!a`, `!!b`, `!!{...}`).
%
% Author text is untouched; only the packaging is ours.
\documentclass[10pt,a4paper]{article}

\usepackage{amsmath}
\usepackage{amssymb}
\usepackage{amsthm}
\usepackage{graphicx}
\usepackage{hyperref}
\usepackage{enumitem}
\usepackage[framemethod=default]{mdframed}

% ---- OLP shorthand accents (from sty/open-logic.sty) ----
% NB: \! is taken by LaTeX (negative thin space), so we only define the !!a / !!b families
% and the \olup helper used by them.
\newcommand*{\olup}[1]{\MakeUppercase{#1}}
\def\!!!a#1{\olup{#1}}
\def\!a#1{\olup{#1}}
\def\!b#1{\olup{#1}}

% ---- theorem-like environments, same family numbering as OLP ----
\theoremstyle{plain}
\newtheorem{thm}{Theorem}[section]
\newtheorem{lem}[thm]{Lemma}
\newtheorem{prop}[thm]{Proposition}
\newtheorem{cor}[thm]{Corollary}
\newtheorem{axiom}{Axiom}

\theoremstyle{definition}
\newtheorem{defn}[thm]{Definition}
\newtheorem{ex}[thm]{Example}
\newtheorem{prob}{Problem}[section]
\newtheorem{probd}{Problem}
\newtheorem{conv}[thm]{Convention}

\theoremstyle{remark}
\newtheorem{rem}[thm]{Remark}
\newtheorem{note}[thm]{Note}
\newtheorem{case}[thm]{Case}

% ---- environments OLP defines as wrappers ----
\newenvironment{probtag}[1]{\begin{prob}}{\end{prob}}
\newenvironment{explain}{}{}
\newenvironment{explain*}{}{}
\newenvironment{editorial}{}{}
\newenvironment{digress}{}{}
\newenvironment{digress*}{}{}
\newenvironment{history}{}{}
\newenvironment{reference}{}{}
\newcommand*{\olfileid}[3]{}
\newcommand*{\olsection}[1]{\section*{#1}}
\newcommand*{\olchapter}[3]{}
\newcommand*{\olpart}[2]{}
\newcommand*{\olimport}[1]{}
\newcommand*{\olinclude}[1]{}
\newcommand*{\olpath}[1]{}
\newcommand*{\ollangid}[1]{}
\newcommand*{\problemtag}[1]{}
\newcommand*{\problemauthor}[1]{}
\newcommand*{\ollabel}[1]{\label{#1}}
% \olref takes 0-3 OPTIONAL arguments plus one mandatory label:
%     \olref{lab}  |  \olref[from]{lab}  |  \olref[from][to]{lab}  |  \olref[a][b][c]{lab}
% A single-argument definition silently swallows "[sth]" and emits a broken reference
% (measured on the host: the text came out as "??iso]x"). The chain below consumes the
% brackets and cross-references only the mandatory argument.
\makeatletter
\newcommand{\olref}{\@ifnextchar[{\olref@i}{\olref@final}}
\def\olref@i[#1]{\@ifnextchar[{\olref@ii}{\olref@final}}
\def\olref@ii[#1]{\@ifnextchar[{\olref@iii}{\olref@final}}
\def\olref@iii[#1]{\olref@final}
\def\olref@final#1{\ref{#1}}
\makeatother
\newcommand*{\olnewpage}{\clearpage}

% ===================== OLP compatibility layer (1) =====================

\makeatletter

% ---------------- proof trees (OLP's prooftree is a bussproofs derivative) ----------------
\providecommand{\AxiomC}[1]{}
\providecommand{\UnaryInfC}[1]{}
\providecommand{\BinaryInfC}[1]{}
\providecommand{\TrinaryInfC}[1]{}
\providecommand{\QuaternaryInfC}[1]{}
\providecommand{\QuinaryInfC}[1]{}
\providecommand{\RightLabel}[1]{}
\providecommand{\LeftLabel}[1]{}
\providecommand{\doubleLine}{}
\providecommand{\noLine}{}
\providecommand{\DisplayProof}{}
\providecommand{\fCenter}{}
\newenvironment{prooftree}{}{}
\newenvironment{ltprooftree}{}{}
\newenvironment{oltableau}{}{}

% ---------------- logical notation ----------------
\providecommand{\lif}{\mathbin{\rightarrow}}
\providecommand{\lforall}{\mathord{\forall}}
\providecommand{\lexists}{\mathord{\exists}}
\providecommand{\ltrue}{\top}
\providecommand{\lfalse}{\bot}
\providecommand{\Ltrue}{\top}
\providecommand{\Lfalse}{\bot}
\providecommand{\Sequent}{\mathrel{\Rightarrow}}
\providecommand{\Proves}{\mathrel{\vdash}}
\providecommand{\Entails}{\mathrel{\models}}
\providecommand{\Subst}[3]{#1[#2/#3]}
\providecommand{\sFmla}[2]{#1(#2)}
\providecommand{\pFmla}[2]{#1(#2)}
\providecommand{\formula}{\relax}
\providecommand{\True}{\top}
\providecommand{\False}{\bot}
\providecommand{\Args}[1]{#1}
\providecommand{\Ts}{\ensuremath{\top}}
\providecommand{\Fi}{\ensuremath{\bot}}
\providecommand{\lfalsum}{\bot}

% ---------------- objects, structures, metatheory ----------------
\providecommand{\Obj}{}
\providecommand{\Struct}{\ensuremath{\mathfrak{A}}}
\providecommand{\Domain}{\ensuremath{|\mathfrak{A}|}}
\providecommand{\Assign}{\ensuremath{s}}
\providecommand{\Sat}[2]{\ensuremath{#1 \models #2}}
\providecommand{\mSat}[3]{\ensuremath{#1 \models #2[#3]}}
\providecommand{\mModel}[2]{\ensuremath{#1 \models #2}}
\providecommand{\Lang}{\ensuremath{\mathcal{L}}}
\providecommand{\Th}{\ensuremath{\mathrm{Th}}}
\providecommand{\Log}{\ensuremath{\mathrm{L}}}
\providecommand{\tuple}[1]{\langle #1 \rangle}
\providecommand{\num}[1]{\underline{#1}}
\providecommand{\cfind}[1]{\varphi_{#1}}
\providecommand{\Nat}{\ensuremath{\mathbb{N}}}
\providecommand{\Mult}{\ensuremath{\mathrm{Mult}}}
\providecommand{\Add}{\ensuremath{\mathrm{Add}}}
\providecommand{\Exp}{\ensuremath{\mathrm{Exp}}}
\providecommand{\Zero}{\ensuremath{0}}
\providecommand{\Succ}{\ensuremath{\mathrm{S}}}
\providecommand{\lambd}{\ensuremath{\lambda}}
\providecommand{\bredpar}{\mathrel{\rightarrow_\beta}}
\providecommand{\bcd}[1]{\ensuremath{\ulcorner #1 \urcorner}}
\providecommand{\Set}[1]{\ensuremath{\{#1\}}}
\providecommand{\indcase}[1]{\par\noindent\emph{Case #1}.\quad}
\providecommand{\indfrm}[1]{\par\noindent\emph{Form #1}.\quad}
\providecommand{\iftag}[1]{#1}
\providecommand{\tagitem}[1]{\item[#1]}
\providecommand{\Elim}[1]{}
\providecommand{\Intro}[1]{}
\providecommand{\LeftR}[1]{}
\providecommand{\RightR}[1]{}
\providecommand{\TAss}[1]{}
\providecommand{\TRule}[1]{}

\makeatother
\providecommand{\olpath}{.}

% ===================== OLP compatibility layer (2) =====================
% --- second compatibility block: notation and shorthands used by the frozen OLP units ---
\providecommand{\CutCS}{\mathrm{Cut}}
\providecommand{\LeftSubproofLabel}{\relax}
\providecommand{\RightSubproofLabel}{\relax}
\providecommand{\maeh}{\relax}
\providecommand{\cheight}{\relax}
\providecommand{\ande}{\mathrm{and}_E}
\providecommand{\andi}{\mathrm{and}_I}
\providecommand{\ore}{\mathrm{or}_E}
\providecommand{\ori}{\mathrm{or}_I}
\providecommand{\typeof}{\mathord{:}}
\providecommand{\pValue}[1]{#1}
\providecommand{\Sent}{\mathrm{Sent}}
\providecommand{\cardplus}{\mathbin{+}}
\providecommand{\cardtimes}{\mathbin{\cdot}}
\providecommand{\cardexpo}{\mathbin{\hat{\ }}}
\providecommand{\cardnless}{\not<}
\providecommand{\cardle}{\leq}
\providecommand{\funfromto}[3]{#1\colon #2 \to #3}
\providecommand{\funimage}[1]{\operatorname{im}(#1)}
\providecommand{\fdefined}{\mathord{\downarrow}}
\providecommand{\canonord}{\operatorname{ord}}
\providecommand{\ordtimes}{\mathbin{\cdot}}
\providecommand{\ordeq}[2]{#1 = #2}
\providecommand{\ordneq}[2]{#1 \neq #2}
\providecommand{\ordtype}[1]{\operatorname{ord}(#1)}
\providecommand{\ordsucc}{\mathord{+1}}
\providecommand{\ordunion}{\cup}
\providecommand{\onesphere}{S^1}
\providecommand{\rotationsgroup}{\mathrm{SO}(2)}
\providecommand{\equivrep}[1]{[#1]}
\providecommand{\rlexless}{\prec}
\providecommand{\ZFC}{\mathrm{ZFC}}
\providecommand{\ZF}{\mathrm{ZF}}
\providecommand{\Rat}{\mathbb{Q}}
\providecommand{\Real}{\mathbb{R}}
\providecommand{\PosInt}{\mathbb{Z}^{+}}
\providecommand{\Id}{\mathrm{id}}
\providecommand{\Bin}{\mathrm{Bin}}
\providecommand{\pto}{\rightharpoonup}
\providecommand{\dom}{\operatorname{dom}}
\providecommand{\ran}{\operatorname{ran}}
\providecommand{\len}{\operatorname{len}}
\providecommand{\concat}{\mathbin{{}^\frown}}
\providecommand{\card}[1]{\lvert #1 \rvert}
\providecommand{\Pow}{\mathcal{P}}
\providecommand{\Setabs}[2]{\{\, #1 \mid #2 \,\}}
\providecommand{\comp}{\circ}
\providecommand{\Char}{\chi}
\providecommand{\Prf}{\operatorname{Prf}}
\providecommand{\Prov}{\operatorname{Prov}}
\providecommand{\Proj}{\operatorname{proj}}
\providecommand{\Gn}{\operatorname{gn}}
\providecommand{\gn}{\operatorname{gn}}
\providecommand{\Value}[1]{\operatorname{val}(#1)}
\providecommand{\Deduce}{\mathrm{D}}
\providecommand{\Dual}{\mathrm{d}}
\providecommand{\Ax}{\mathrm{Ax}}
\providecommand{\pSat}[3]{#1 \models #2\,[#3]}
\providecommand{\pAssign}{\sigma}
\providecommand{\bforall}{\forall}
\providecommand{\bexists}{\exists}
\providecommand{\liff}{\mathbin{\leftrightarrow}}
\providecommand{\FV}{\operatorname{FV}}
\providecommand{\UnaryInf}{\ensuremath{\mid}}
\providecommand{\BinaryInf}{\ensuremath{\mid}}
\providecommand{\depth}{\operatorname{depth}}
\providecommand{\oliflabeldef}[3]{#1}
\providecommand{\TMblank}{\sqcup}
\providecommand{\TMstroke}{\mid}
\providecommand{\TMstay}{\mathrm{S}}
\providecommand{\TMright}{\mathrm{R}}
\providecommand{\TMleft}{\mathrm{L}}
\providecommand{\TMendtape}{\triangleright}
\providecommand{\umin}{\mu}
\providecommand{\Sum}{\sum}
\providecommand{\citep}[1]{\cite{#1}}
\providecommand{\citet}[1]{\cite{#1}}
\providecommand{\citeauthor}[1]{\cite{#1}}
\providecommand{\citeyear}[1]{\cite{#1}}

% ===================== OLP compatibility layer (3) =====================
% --- third compatibility block: remaining notation seen in the frozen OLP units ---
\providecommand{\iso}{\cong}
\providecommand{\nszero}{\mathrm{0}}
\providecommand{\nssucc}{\mathrm{S}}
\providecommand{\nsplus}{\mathbin{+}}
\providecommand{\nssless}{<}
\providecommand{\nsless}{<}
\providecommand{\PL}{\mathrm{PL}}
\providecommand{\Nec}{\mathord{\Box}}
\providecommand{\RK}{\mathrm{RK}}
\providecommand{\ST}{\mathrm{ST}}
\providecommand{\SSubst}[3]{#1[#2/#3]}
\providecommand{\subst}[3]{#1[#2/#3]}
\providecommand{\tagfalse}{\mathrm{F}}
\providecommand{\cutr}{\mathrm{Cut}}
\providecommand{\cardneq}{\neq}
\providecommand{\issucc}[1]{\operatorname{succ}(#1)}
\providecommand{\tup}[1]{\langle #1 \rangle}
\providecommand{\Var}{\mathrm{Var}}
\providecommand{\Bew}{\operatorname{Bew}}
\providecommand{\Con}{\operatorname{Con}}
\providecommand{\Cons}{\operatorname{Cons}}
\providecommand{\pDeduce}{\mathrm{D}}
\providecommand{\psat}{\models}
\providecommand{\den}[1]{\llbracket #1 \rrbracket}
\providecommand{\Sem}[2]{\llbracket #1 \rrbracket_{#2}}
\providecommand{\DomainOf}[1]{\lvert #1 \rvert}
\providecommand{\Form}{\mathrm{Form}}
\providecommand{\Terms}{\mathrm{Term}}
\providecommand{\Atom}{\mathrm{Atom}}
\providecommand{\FVars}[1]{\operatorname{FV}(#1)}
\providecommand{\fv}[1]{\operatorname{fv}(#1)}
\providecommand{\tuplep}[1]{\langle #1 \rangle}
\providecommand{\restrict}[2]{#1 \upharpoonright #2}
\providecommand{\sto}{\rightarrow}
\providecommand{\FmlaSet}[1]{\mathrm{Fmla}_{#1}}
\providecommand{\Thm}{\mathrm{Thm}}
\providecommand{\Eval}{\operatorname{val}}
\providecommand{\blank}{\sqcup}
\providecommand{\blankTape}{\sqcup}
\providecommand{\Consis}{\operatorname{Con}}
\providecommand{\drs}{\mathrm{d}}
\providecommand{\VRule}{\mid}
\providecommand{\Bet}{\mathrm{B}}
\providecommand{\Seg}[1]{\mathrm{Seg}_{#1}}
\providecommand{\Enums}{\mathrm{En}}
\providecommand{\Z}{\mathbb{Z}}
\providecommand{\N}{\mathbb{N}}
\providecommand{\abs}[1]{\lvert #1 \rvert}
\providecommand{\floor}[1]{\lfloor #1 \rfloor}
\providecommand{\ceil}[1]{\lceil #1 \rceil}

% ===================== OLP compatibility layer (4) =====================
% --- fourth compatibility block: harvested from the per-candidate compile logs ---
\providecommand{\beredpar}{\mathrel{\rightarrow_\beta}}
\makeatletter
\def\bered{\@ifnextchar[{\bered@i}{\mathrel{\rightarrow_\beta}}}
\def\bered@i[#1]{\mathrel{\rightarrow_\beta}}
\def\bred{\@ifnextchar[{\bred@i}{\mathrel{\rightarrow_\beta}}}
\def\bred@i[#1]{\mathrel{\rightarrow_\beta}}
\makeatother
\providecommand{\beredone}{\mathrel{=_\beta}}
\providecommand{\bredone}{\mathrel{=_\beta}}
\providecommand{\xredone}{\mathrel{=_\beta}}
\providecommand{\aconvone}{\mathrel{=_\alpha}}
\providecommand{\cardless}{<}
\providecommand{\tagrefs}[1]{#1}
\providecommand{\ordplus}{\mathbin{+}}
\providecommand{\funrestrictionto}[2]{#1 \upharpoonright #2}
\providecommand{\fundefined}{\mathord{\uparrow}}
\providecommand{\Trm}{\mathrm{Trm}}
\providecommand{\OProv}{\operatorname{Prov}}
\providecommand{\OPrf}{\operatorname{Prf}}
\providecommand{\aeq}{\mathrel{=}}
\providecommand{\mClass}[1]{[#1]}
\providecommand{\equivclass}[1]{[#1]}
\providecommand{\pheight}{\operatorname{ht}}
\providecommand{\phy}{\operatorname{ht}}
\providecommand{\Cut}{\mathrm{Cut}}
\providecommand{\QuantRank}{\operatorname{qr}}
\providecommand{\elemequiv}{\mathrel{\equiv}}
\providecommand{\PIso}{\mathrm{PIso}}
\providecommand{\Theory}{\mathrm{Th}}
\providecommand{\supstrict}{\subsetneq}
\providecommand{\trcl}{\operatorname{trcl}}
\providecommand{\cref}[1]{\ref{#1}}
\providecommand{\Discharge}{\mathrm{Dis}}
\providecommand{\Weakening}{\mathrm{Wk}}
\providecommand{\ZFminus}{\mathrm{ZF}^{-}}
\makeatletter
\renewcommand{\Dual}{\@ifnextchar[{\Dual@i}{\ifmmode\mathrm{d}\else$\mathrm{d}$\fi}}
\def\Dual@i[#1]{\ifmmode\mathrm{d}[#1]\else$\mathrm{d}[#1]$\fi}
\makeatother
\providecommand{\RightR}{\mathrm{R}}
\providecommand{\LeftR}{\mathrm{L}}
\providecommand{\Intro}{\mathrm{I}}
\providecommand{\Axiom}{\mathrm{Ax}}
\providecommand{\formula}[1]{#1}
\providecommand{\stared}{\ast}
\providecommand{\oset}[2]{\{\, #1 \mid #2 \,\}}
\providecommand{\setof}[2]{\{\, #1 \mid #2 \,\}}
\providecommand{\pow}[1]{\mathcal{P}(#1)}
\providecommand{\ordpair}[2]{\langle #1, #2 \rangle}
\providecommand{\ordtriple}[3]{\langle #1, #2, #3 \rangle}
\providecommand{\union}{\cup}
\providecommand{\inter}{\cap}
\providecommand{\seq}[1]{\langle #1 \rangle}
\providecommand{\SuchThat}{\mid}
\providecommand{\where}{\mid}
\providecommand{\TVar}{\mathrm{TVar}}
\providecommand{\Bool}{\mathrm{Bool}}
\providecommand{\Val}{\operatorname{Val}}
\providecommand{\TL}[1]{\mathrm{TL}}
\providecommand{\xstared}{\ast}

% ===================== OLP compatibility layer (5) =====================
% --- fifth compatibility block: rescues the quarantined OLP candidates ---
\providecommand{\ident}{\operatorname{id}}
\providecommand{\setrank}{\operatorname{rank}}
\providecommand{\cardfont}{\mathfrak{c}}
\providecommand{\Complement}{\overline}
\providecommand{\QR}{\mathrm{QR}}
\providecommand{\becd}{\ensuremath{\ulcorner}}
\providecommand{\disjointsum}{\mathbin{\oplus}}
\providecommand{\cardeq}[2]{#1 = #2}
\providecommand{\xred}{\mathrel{\rightarrow}}
% --- sixth compatibility block: names reported by real compile logs ---
\providecommand{\ORefut}{\operatorname{Refut}}
\providecommand{\ORProv}{\operatorname{Prov}}
\providecommand{\cutrank}{\operatorname{cr}}
\providecommand{\Refut}{\operatorname{Refut}}
\providecommand{\DischargeRule}{\mathrm{Dis}}
\providecommand{\Contraction}{\mathrm{C}}
\providecommand{\bottomAlignProof}{\relax}
\providecommand{\FalseCl}{\bot}
\providecommand{\red}{\mathrel{\rightarrow}}
\providecommand{\varAssign}{\sigma}
\providecommand{\defiff}{\mathrel{\leftrightarrow}}
\providecommand{\FalseInt}{\bot}
\providecommand{\fn}{\operatorname{fn}}
\providecommand{\dotted}{\cdot}

% --- tag-macro semantic fix (2026-10-13, final): \iftag CHOOSES a branch and the OLP proof body is
% always the SECOND group (\iftag{tag}{body}, {tag}{body}{c}, {tag}{body}{c}{d}); a definition that
% prints the tag literal drops the proof body. Bodies contain blank lines, hence \long.
\makeatletter
\providecommand{\iftag}{}
\long\def\iftag#1#2{\@ifnextchar\bgroup{\iftag@three{#2}}{\iftag@stop{#2}}}
\long\def\iftag@three#1#2{\iftag@stop{#1}\@ifnextchar\bgroup{\iftag@eat}{\relax}}
\long\def\iftag@eat#1{}
\long\def\iftag@stop#1{#1}
\long\def\oliflabeldef#1#2{\@ifnextchar\bgroup{\@gobble}{\iftag@stop{#2}}}
\providecommand{\tagrefs}[1]{}
\providecommand{\settag}[1]{}
\makeatother
\begin{document}

% ==== OLP file: content/second-order-logic/metatheory/compactness.tex ====
% Part: second-order-logic
% Chapter: metatheory
% Section: compactness



\olfileid{sol}{met}{com}

\olsection{Second-order Logic is not Compact}

\begin{explain}
Call a set of sentences~$\Gamma$ \emph{finitely satisfiable} if every
one of its finite subsets is satisfiable.  First-order logic has the
property that if a set of \olup{sentence}s~$\Gamma$ is finitely
satisfiable, it is satisfiable.  This property is called
\emph{compactness}.  It has an equivalent version involving
entailment: if $\Gamma \Entails A$, then already $\Gamma_0 \Entails
A$ for some finite subset~$\Gamma_0 \subseteq \Gamma$. In this
version it is an immediate corollary of the completeness theorem: for
if $\Gamma \Entails A$, by completeness $\Gamma \Proves A$. But a
\olup{derivation} can only make use of finitely many \olup{sentence}s
of~$\Gamma$.

Compactness is not true for second-order logic. There are sets of
second-order \olup{sentence}s that are finitely satisfiable but not
satisfiable, and that entail some~$A$ without a finite subset
entailing~$A$.  
\end{explain}


\begin{thm}
\ollabel{thm:sol-undecidable}
Second-order logic is not compact.
\end{thm}

\begin{proof}
Recall that 
\[
\fn{Inf} \ident \lexists[u][(\lforall[x][\lforall[y][(\eq[u(x)][u(y)]
      \lif \eq[x][y])]] \land \lexists[y][\lforall[x][\eq/[y][u(x)]]])]
\]
is satisfied in \olup{structure} iff its domain is infinite. Let
$A^{\ge n}$ be a sentence that asserts that the domain has at least
$n$ \olup{element}s, e.g.,
\[
A^{\ge n} \ident \lexists[x_1][\dots\lexists[x_n][(\eq/[x_1][x_2]
    \land \eq/[x_1][x_3] \land \dots \land \eq/[x_{n-1}][x_n])]].
\]
Consider the set of \olup{sentence}s
\[
\Gamma = \{\lnot \fn{Inf}, A^{\ge 1}, A^{\ge 2}, A^{\ge 3}, \dots\}.
\]
It is finitely satisfiable, since for any finite subset~$\Gamma_0
\subseteq \Gamma$
there is some $k$ so that $A^{\ge k} \in \Gamma$ but no $A^{\ge n}
\in \Gamma$ for $n>k$. If $\Domain{M}$ has $k$ \olup{element}s,
$\Sat{M}{\Gamma_0}$. But, $\Gamma$ is not satisfiable: if
$\Sat{M}{\lnot \fn{Inf}}$, $\Domain{M}$ must be finite, say, of
size~$k$. Then $\Sat/{M}{A^{\ge k+1}}$.
\end{proof}

\begin{prob}
Give an example of a set~$\Gamma$ and \olup{sentence}~$A$ so that
$\Gamma \Entails A$ but for every finite subset~$\Gamma_0 \subseteq
\Gamma$, $\Gamma_0 \Entails/ A$.
\end{prob}



% ==== OLP file: content/second-order-logic/metatheory/introduction.tex ====
% Part: second-order-logic
% Chapter: metatheory
% Section: introduction



\olfileid{sol}{met}{int}

\olsection{Introduction}

First-order logic has a number of nice properties. We know it is not
decidable, but at least it is axiomatizable. That is, there are proof
systems for first-order logic which are sound and complete, i.e., they
give rise to \olup{derivability} relation~$\Proves$ with the property
that for any set of \olup{sentence}s~$\Gamma$ and \olup{sentence}~$A$,
$\Gamma \Entails A$ iff $\Gamma \Proves A$.  This means in
particular that the validities of first-order logic are \olup{computably
  enumerable}. There is a computable function~$f\colon \Nat \to
\Sent[L]$ such that the values of~$f$ are all and only the valid
\olup{sentence}s of~$\Lang{L}$. This is so because \olup{derivation}s can be
enumerated, and those that \olup{derive} a single~\olup{sentence} are then
mapped to that \olup{sentence}.  Second-order logic is more expressive
than first-order logic, and so it is in general more complicated to
capture its validities.  In fact, we'll show that second-order logic
is not only undecidable, but its validities are not even \olup{computably
  enumerable}. This means there can be no sound and complete proof
system for second-order logic (although sound, but incomplete proof
systems are available and in fact are important objects of research).

First-order logic also has two more properties: it is compact (if
every finite subset of a set~$\Gamma$ of \olup{sentence}s is satisfiable,
$\Gamma$ itself is satisfiable) and the L\"owenheim--Skolem Theorem
holds for it (if $\Gamma$ has an infinite model it has
\olup{denumerable} model). Both of these results fail for second-order
logic. Again, the reason is that second-order logic can express facts
about the size of \olup{domain}s that first-order logic cannot.


% ==== OLP file: content/second-order-logic/metatheory/loewenheim-skolem.tex ====
% Part: second-order-logic
% Chapter: metatheory
% Section: compactness



\olfileid{sol}{met}{lst}

\olsection{The L\"owenheim--Skolem Theorem Fails for Second-order Logic}

\begin{explain}
The (Downward) L\"owenheim--Skolem Theorem states that every set of
\olup{sentence}s with an infinite model has \olup{enumerable} model.  It,
too, is a consequence of the completeness theorem: the proof of
completeness generates a model for any consistent set of
\olup{sentence}s, and that model is \olup{enumerable}.  There is also an
Upward L\"owenheim--Skolem Theorem, which guarantees that if a set of
\olup{sentence}s has \olup{denumerable} model it also has
\olup{nonenumerable} model.  Both theorems fail in second-order logic.
\end{explain}


\begin{thm}
\ollabel{thm:sol-no-ls} The L\"owenheim--Skolem Theorem fails for
second-order logic: There are \olup{sentence}s with infinite models but
no \olup{enumerable} models.
\end{thm}

\begin{proof}
Recall that 
\[
\fn{Count} \ident \lexists[z][\lexists[u][\lforall[X][((X(z) \land
      \lforall[x][(X(x) \lif X(u(x)))]) \lif \lforall[x][X(x)])]]]
\]
is true in \olup{structure}~$\Struct{M}$ iff $\Domain{M}$ is
\olup{enumerable}, so $\lnot\fn{Count}$ is true in~$\Struct{M}$
iff $\Domain{M}$ is \olup{nonenumerable}.  There are
such \olup{structure}s---take any \olup{nonenumerable} set as the
\olup{domain}, e.g., $\Pow{\Nat}$ or~$\Real$. So $\lnot \fn{Count}$ 
has infinite models but no \olup{enumerable} models.
\end{proof}

\begin{thm}
There are \olup{sentence}s with \olup{denumerable} but no
\olup{nonenumerable} models.
\end{thm}

\begin{proof}
$\fn{Count} \land \fn{Inf}$ is true in $\Nat$ but not in any
  \olup{structure}~$\Struct{M}$ with $\Domain{M}$ \olup{nonenumerable}.
\end{proof}



% ==== OLP file: content/second-order-logic/metatheory/metatheory.tex ====
% Part: second-order-logic
% Chapter: metatheory



\olchapter{sol}{met}{Metatheory of Second-order Logic}

\olimport{introduction}

\olimport{second-order-arithmetic}

\olimport{undecidability-and-axiomatizability}

\olimport{compactness}

\olimport{loewenheim-skolem}

\OLEndChapterHook


% ==== OLP file: content/second-order-logic/metatheory/second-order-arithmetic.tex ====
% Part: second-order-logic
% Chapter: metatheory
% Section: second-order-arithmetic



\olfileid{sol}{met}{spa}

\olsection{Second-order Arithmetic}

Recall that the theory $\Th{PA}$ of Peano arithmetic includes the
eight axioms of~$\Th{Q}$,
\begin{align*}
& \lforall[x][\eq/[x'][\Obj 0]]\\
& \lforall[x][\lforall[y][(\eq[x'][y'] \lif \eq[x][y])]]\\
& \lforall[x][(\eq[x][\Obj 0] \lor \lexists[y][\eq[x][y']])]\\ 
& \lforall[x][\eq[(x + \Obj 0)][x]]\\
& \lforall[x][\lforall[y][\eq[(x + y')][(x + y)']]]\\
& \lforall[x][\eq[(x \times \Obj 0)][\Obj 0]]\\
& \lforall[x][\lforall[y][\eq[(x \times y')][((x \times y) + x)]]]\\
& \lforall[x][\lforall[y][(x < y \liff \lexists[z][\eq[(z' + x)][y]])]]\\
\intertext{plus all sentences of the form}
& (A(\Obj 0) \land \lforall[x][(A(x) \lif A(x'))]) \lif \lforall[x][A(x)].
\end{align*}
The latter is a ``schema,'' i.e., a pattern that generates infinitely
many \olup{sentence}s of the language of arithmetic, one for each
\olup{formula}~$A(x)$. We call this schema the (first-order)
\emph{axiom schema of induction}. 
In \emph{second-order} Peano arithmetic~$\Th{PA^2}$, induction
can be stated as a single sentence. $\Th{PA^2}$
consists of the first eight axioms above plus the (second-order)
\emph{induction axiom}:
\[
\lforall[X][((X(\Obj 0) \land \lforall[x][(X(x) \lif X(x'))]) \lif \lforall[x][X(x)])].
\]
It says that if a subset~$X$ of the \olup{domain}
contains~$\Assign{\Obj{0}}{M}$ and with any~$x \in \Domain{M}$ also
contains~$\Assign{\prime}{M}(x)$ (i.e., it is ``closed under
successor'') it contains everything in the \olup{domain} (i.e., $X =
\Domain{M})$.

The induction axiom guarantees that any \olup{structure} satisfying it
contains only those \olup{element}s of~$\Domain{M}$ the axioms require to
be there, i.e., the values of~$\num{n}$ for $n \in \Nat$. A model of
$\Th{PA^2}$ contains no non-standard numbers.

\begin{thm}
\ollabel{thm:sol-pa-standard}
If $\Sat{M}{\Th{PA^2}}$ then $\Domain{M} =
\Setabs{\Value{\num{n}}{M}}{n \in \Nat}$.
\end{thm}

\begin{proof}
Let $N = \Setabs{\Value{\num{n}}{M}}{n \in \Nat}$, and suppose
$\Sat{M}{\Th{PA^2}}$. Of course, for any $n \in \Nat$,
$\Value{\num{n}}{M} \in \Domain{M}$, so $N \subseteq \Domain{M}$.

Now for inclusion in the other direction. Consider a variable
assignment $s$ with $s(X) = N$. By assumption,
\begin{align*}
\Sat{M}{\lforall[X][((X(\Obj 0) \land \lforall[x][(X(x)
      \lif X(x'))]) \lif \lforall[x][X(x)])]}, & \text{ thus}\\
\Sat{M}{(X(\Obj 0) \land \lforall[x][(X(x) \lif X(x'))]) \lif
  \lforall[x][X(x)]}[s]. & 
\end{align*}
Consider the antecedent of this conditional.  $\Value{\Obj{0}}{M} \in
N$, and so $\Sat{M}{X(\Obj{0})}[s]$.  The second conjunct,
$\lforall[x][(X(x) \lif X(x'))]$ is also satisfied. For suppose $x \in
N$. By definition of~$N$, $x = \Value{\num{n}}{M}$ for some~$n$.  That
gives $\Assign{\prime}{M}(x) = \Value{\num{n+1}}{M} \in N$. So,
$\Assign{\prime}{M}(x) \in N$.

We have that $\Sat{M}{X(\Obj 0) \land \lforall[x][(X(x) \lif X(x'))]}[s]$.
Consequently, $\Sat{M}{\lforall[x][X(x)]}[s]$. But that means that for
every $x \in \Domain{M}$ we have $x \in s(X) = N$. So, $\Domain{M}
\subseteq N$.
\end{proof}

\begin{cor}
\ollabel{cor:sol-pa-categorical}
Any two models of~$\Th{PA^2}$ are isomorphic.
\end{cor}

\begin{proof}
By \olref{thm:sol-pa-standard}, the domain of any model of $\Th{PA^2}$
is exhausted by $\Value{\num{n}}{M}$. Any such model is also a model
of~$\Th{Q}$. By \olref[mod][mar][stm]{prop:thq-standard}, any such
model is standard, i.e., isomorphic to~$\Struct{N}$.
\end{proof}

Above we defined $\Th{PA^2}$ as the theory that contains the first eight 
arithmetical axioms plus the second-order induction axiom.  In fact, thanks
to the expressive power of second-order logic, only the \emph{first two} 
of the arithmetical axioms plus induction are needed for second-order Peano 
arithmetic.

\begin{prop}\ollabel{prop:sol-pa-definable}
Let $\Th{PA^{2\dagger}}$ be the second-order theory containing the
first two arithmetical axioms (the successor axioms) and the
second-order induction axiom. Then $\le$, $+$, and $\times$ are definable
in $\Th{PA^{2\dagger}}$.
\end{prop}

\begin{proof}
To show that $\le$ is definable, we have to find a formula~$A_\le(x, y)$
such that $\Sat{N}{A_\le(\num n, \num m)}$ iff $n \le m$. Consider the formula
\[
B(x, Y) \ident Y(x) \land \lforall[y][(Y(y) \lif Y(y'))]
\]
Clearly, $B(\num n, Y)$ is satisfied by a set $Y \subseteq \Nat$ iff
$\Setabs{m}{n \le m} \subseteq Y$, so we can take $A_\le(x, y) \ident
\lforall[Y][(B(x, Y) \lif Y(y))]$.

To see that addition is definable observe that $k+l=m$ iff there is a
function $u$ such that $u(0)=k$, $u(n')=u(n)'$ for all $n$, and $m = u(l)$. 
We can use this equivalence to define addition in $\Th{PA^{2\dagger}}$ 
by the following formula:
\[
A_+(x,y,z) \ident \lexists[u][(u(\Obj 0)=x \land \lforall[w][u(x')=u
 (x)'] \land u(y)=z)]
\] 
It should be clear that $\Sat{N}{A_+(\num k, \num l, \num m)}$ 
iff $k+l=m$.
\end{proof}

\begin{prob}
Complete the proof of \olref[sol][met][spa]{prop:sol-pa-definable}.
\end{prob}


% ==== OLP file: content/second-order-logic/metatheory/undecidability-and-axiomatizability.tex ====
% Part: second-order-logic
% Chapter: metatheory
% Section: undecidability-and-axiomatizability



\olfileid{sol}{met}{nax}

\olsection{Second-order Logic is not Axiomatizable}


\begin{thm}
\ollabel{thm:sol-undecidable}
Second-order logic is undecidable.
\end{thm}

\begin{proof}
A first-order \olup{sentence} is valid in first-order logic iff it is
valid in second-order logic, and first-order logic is undecidable.
\end{proof}

\begin{thm}
\ollabel{cor:sol-not-axiomatizable}
There is no sound and complete \olup{derivation} system for second-order logic.
\end{thm}

\begin{proof}
Let $A$ be \olup{sentence} in the language of arithmetic.  $\Sat{N}{A}$
iff $\Th{PA^2} \Entails A$. Let $P$ be the conjunction of the nine
axioms of~$\Th{PA^2}$. $\Th{PA^2} \Entails A$ iff $\Entails P \lif
A$, i.e., $\Sat{M}{P \lif A$}. Now consider the \olup{sentence}
$\lforall[z][\lforall[u][\lforall[u'][\lforall[u''][\lforall[L][(P' \lif A')]]]]]$ resulting by
replacing $\Obj{0}$ by $z$, $\prime$ by the one-place function
variable $u$, $+$ and $\times$ by the two-place function-variables 
$u'$ and  $u''$, respectively,
and $<$ by the two-place relation variable~$L$ and
universally quantifying.  It is a valid sentence of pure second-order
logic iff the original sentence was valid iff $\Th{PA^2} \Entails A$
iff $\Sat{N}{A}$.  Thus if there were a sound and complete proof
system for second-order logic, we could use it to define a computable
enumeration $f\colon \Nat \to \Sent[L_A]$ of the \olup{sentence}s true
in~$\Struct{N}$. This function would be representable in~$\Th{Q}$ by
some first-order formula~$B_f(x, y)$. Then the
\olup{formula}~$\lexists[x][B_f(x, y)]$ would define the set of true
first-order \olup{sentence}s of~$\Struct{N}$, contradicting Tarski's
Theorem.
\end{proof}




\end{document}
